\chapter{\textit{Coherence1.0} typing rules} 
\section{Typing context} \label{appendix:typing_context}
The term \textit{metadata structure} of an object refers to a structure containing all the necessary data about the object. The underlying mathematical representation of these structures are omitted for brevity. $\Gamma$ contains the following:
\begin{itemize}
    \item $\Gamma_V$: a mapping from local variable names to their types
    \item $\Gamma_L$: boolean flag indicating whether the current typing context is scoped within an \texttt{atomic} section
    \item $\Gamma_A$: a mapping from actor names to their metadata structure.
    \item $\Gamma_T$: mapping to named types.
    \item $\Gamma_F$: mapping from top-level function names to their metadata structure.
    \item $\Gamma_{\text{this}}$: an optional metadata structure of the current actor.
    \item $\Gamma_{callable}$: a metadata structure of the current callable (function, constructor, or behavior) being type-checked
\end{itemize}
\section{Expression typing} \label{appendix:expression_typing}
\begin{tabular}{@{\qquad\qquad}c@{\qquad}c@{\qquad}c@{\qquad}c} 
    \AxiomC{}
    \RightLabel{(Int)}
    \UnaryInfC{$\Gamma \tack i : \text{int}$}
    \DisplayProof
&
    \AxiomC{}
    \RightLabel{(True)}
    \UnaryInfC{$\Gamma \tack \text{true} : \text{bool}$}
    \DisplayProof
&
    \AxiomC{}
    \RightLabel{(False)}
    \UnaryInfC{$\Gamma \tack \text{false} : \text{bool}$}
    \DisplayProof
\end{tabular}

\vspace{1em}

\begin{prooftree}
    \AxiomC{$\text{variable\_type}(x, \Gamma) = T$}
    \RightLabel{(Var)}
    \UnaryInfC{$\Gamma \tack x : T$}
\end{prooftree}

\vspace{1em}
\begin{prooftree}
    \AxiomC{$\text{get\_struct\_type}(N, \Gamma) = \text{struct } \{ \overline{f:B} \}$}
    \AxiomC{$\Gamma \tack \overline{e} : \overline{B'}$}
    \AxiomC{$\Gamma \tack \text{assignable}(\overline{B}, \overline{B'})$}
    \RightLabel{(Struct)}
    \TrinaryInfC{$\Gamma \tack (\{ \overline{f = e} \} : N) : N$}
\end{prooftree}
\vspace{1em}

\begin{prooftree}
    \AxiomC{$\Gamma \tack e_s : \text{Int}$}
    \AxiomC{$\Gamma \tack e_d : B'$}
    \AxiomC{$\Gamma \tack \text{assignable}(B, B')$}
    \RightLabel{(Alloc)}
    \TrinaryInfC{$\Gamma \tack \text{new } \kappa[e_s] B(e_d) : \text{unalias}(B \ \kappa)$}
\end{prooftree}
\vspace{1em}

\begin{prooftree}
    \AxiomC{$A \in \Gamma_A$}
    \AxiomC{$k \in A_{\text{ctor}}$}
    \AxiomC{$\Gamma \tack \overline{e} : \overline{T_{\text{arg}}}$}
    \AxiomC{$\Gamma \tack \text{assignable}(k_{\text{param}}, \overline{T_{\text{arg}}})$}
    \RightLabel{(Actor creation)}
    \QuaternaryInfC{$\Gamma \tack \text{new } A.k(\overline{e}) : \text{Actor}(A)$}
\end{prooftree}

\vspace{1em}
\begin{prooftree}
    \AxiomC{$\Gamma \tack e_p : T\ \kappa$}
    \AxiomC{$\Gamma \tack e_i : \text{Int}$}
    % \AxiomC{$\Gamma_L = 1$}
    \AxiomC{$\Gamma \tack \text{cap\_readable}(\kappa)$}
    \RightLabel{(Pointer access)}
    \TrinaryInfC{$\Gamma \tack e_p[e_i] : T$}
\end{prooftree}

\vspace{1em}

\begin{prooftree}
    \AxiomC{$\Gamma \tack e : N$}
    \AxiomC{$S(N) = \text{struct } \{ \overline{f:B} \}$}
    \AxiomC{$f_i \in \overline{f}$}
    \RightLabel{(Field access)}
    \TrinaryInfC{$\Gamma \tack e.f_i : B_i$}
\end{prooftree}

\vspace{1em}

\begin{prooftree}
    \AxiomC{$\Gamma \tack e_1 : T_1$}
    \AxiomC{$\Gamma \tack e_2 : T_2$}
    \AxiomC{$\Gamma \tack \text{assignable}(T_1, T_2)$}
    \AxiomC{$\Gamma \tack \text{appear\_lhs}(e_1)$}
    \RightLabel{(Assign)}
    \QuaternaryInfC{$\Gamma \tack e_1 = e_2 : \text{unalias}(T_1)$}
\end{prooftree}

\vspace{1em}

\begin{prooftree}
    \AxiomC{$\texttt{get\_func}(\Gamma, m) = (\overline{T_{\text{param}}} \Rightarrow T_{\text{ret}})$}
    \AxiomC{$\Gamma \tack \overline{e} : \overline{T_{\text{arg}}}$}
    \AxiomC{$\Gamma \tack \text{assignable}(\overline{T_{\text{param}}}, \overline{T_{\text{arg}}})$}
    \RightLabel{(Func-Call)}
    \TrinaryInfC{$\Gamma \tack m(\overline{e}) : T_{\text{ret}} $}
\end{prooftree}

\begin{prooftree}
    \AxiomC{$\Gamma_V(x) = T$}
    \RightLabel{(Unalias)}
    \UnaryInfC{$\Gamma \tack \mathtt{unalias}(x) : \text{unalias}(T)$}
\end{prooftree}

\section{Statement typing} \label{appendix:statement_typing}
\begin{center}
\begin{tabular}{c@{\qquad\qquad}c}
    \AxiomC{}
    \RightLabel{(Seq-Nil)}
    \UnaryInfC{$\Gamma \tack [] \tackl \Gamma$}
    \DisplayProof
&
    \AxiomC{$\Gamma \tack s \tackl \Gamma'$}
    \AxiomC{$\Gamma' \tack \overline{s} \tackl \Gamma''$}
    \RightLabel{(Seq-Cons)}
    \BinaryInfC{$\Gamma \tack s{::}\overline{s} \tackl \Gamma''$}
    \DisplayProof
\end{tabular}
\end{center}
\begin{center}
\begin{prooftree}
    \AxiomC{$\Gamma \tack e : T_e$}
    \AxiomC{$\Gamma \tack \text{assignable}(T, T_e)$}
    \RightLabel{(Var-Decl)}
    \BinaryInfC{$\Gamma \tack \text{var } x : T = e; \tackl \Gamma[\Gamma_V \mapsto \Gamma_V[x \mapsto T]]$}
\end{prooftree}


\begin{prooftree}
    \AxiomC{$\text{is\_constructor}(\Gamma_{callable})$}
    \AxiomC{$(\Gamma_{this})_{\text{fields}}(f) = T_f$}
    \AxiomC{$\Gamma \tack e : T_e$}
    \AxiomC{$\Gamma \tack \text{assignable}(T_f, T_e)$}
    \RightLabel{(Member-Init)}
    \QuaternaryInfC{$\Gamma \tack f := e \tackl \Gamma$}
\end{prooftree}

\begin{prooftree}
    \AxiomC{$\Gamma \tack e_a : \text{Actor}(A_{\text{name}})$}
    \AxiomC{$b \in \Gamma_A(A_{\text{name}})_{behv}$}
    \AxiomC{$\Gamma \tack \overline{e} : \overline{T_e}$}
    \AxiomC{$\Gamma \tack \text{assignable}(b_{\text{param}}, \overline{T_e})$}
    \RightLabel{[Behv-Call]}
    \QuaternaryInfC{$\Gamma \tack e_a \to b(\overline{e}); \tackl \Gamma$}
\end{prooftree}


\begin{tabular}{c@{\qquad\qquad}c}
    \AxiomC{$\Gamma \tack e : \text{Int}$}
    \RightLabel{(Out)}
    \UnaryInfC{$\Gamma \tack \text{OUT } e; \tackl \Gamma$}
    \DisplayProof
&
    \AxiomC{$\Gamma \tack e : T$}
    \RightLabel{[Raw expression]}
    \UnaryInfC{$\Gamma \tack e; \tackl \Gamma$}
    \DisplayProof
\end{tabular}

\begin{prooftree}
    \AxiomC{$\Gamma \tack e : \text{bool}$}
    \AxiomC{$\Gamma \tack \overline{s_1} \tackl \Gamma_1$}
    \AxiomC{$\Gamma \tack \overline{s_2} \tackl \Gamma_2$}
    \RightLabel{(If-Else)}
    \TrinaryInfC{$\Gamma \tack \mathtt{if}\ (e)\ \{\ \overline{s_1}\ \}\ \mathtt{else}\ \{\ \overline{s_2}\ \} \tackl \Gamma$}
\end{prooftree}


\begin{prooftree}
    \AxiomC{$\Gamma[L \mapsto \text{true}] \tack \overline{s} \tackl \Gamma'$}
    \RightLabel{[Atomic]}
    \UnaryInfC{$\Gamma \tack \text{atomic } \{ \overline{s} \} \tackl \Gamma$}
\end{prooftree}


\begin{prooftree}
    \AxiomC{$\Gamma \tack e : \text{Bool}$}
    \AxiomC{$\Gamma \tack \overline{s} \tackl \Gamma'$}
    \RightLabel{[While]}
    \BinaryInfC{$\Gamma \tack \text{while } (e) \{ \overline{s} \} \tackl \Gamma$}
\end{prooftree}


\begin{prooftree}
    \AxiomC{$\Gamma_{callable} = (\overline{T_{\text{param}}} \Rightarrow T_{\text{ret}})$}
    \AxiomC{$\Gamma \tack e : T_{\text{ret}}$}
    \RightLabel{[Return]}
    \BinaryInfC{$\Gamma \tack \text{return } e; \tackl \Gamma$}
\end{prooftree}
\end{center}

\section{Top-level typing} \label{appendix:top_level_item_typing}

\begin{center}
\begin{prooftree}
    \AxiomC{$\Gamma \tack \text{valid}(T_N)$}
    \RightLabel{[TypeDef]}
    \UnaryInfC{$\Gamma \tack \text{type } N = T_N \tack \Gamma[C \mapsto C[N \mapsto T_N]]$}
\end{prooftree}

\vspace{1em}

\begin{prooftree}
    \AxiomC{$\Gamma' = \Gamma[V \mapsto \{\overline{x : \phi_p}\}]$}
    \AxiomC{$\Gamma' \tack \text{returns}(\overline{s}, \phi_r)$}
    \AxiomC{$\Gamma' \tack \text{valid\_var\_usage}(\overline{s})$}
    \AxiomC{$\Gamma' \tack \overline{s} \tackl \Gamma''$}
    \RightLabel{[T-Func]}
    \QuaternaryInfC{$\Gamma \tack \text{func } m(\overline{x : \phi_p}) \Rightarrow \phi_r \{ \overline{s} \} \tackl \Gamma$}
\end{prooftree}

\vspace{1em}

\begin{prooftree}
    \AxiomC{$\Gamma' = \Gamma[V \mapsto \{\overline{x : \phi_p}\}]$}
    \AxiomC{$\Gamma' \tack \text{valid\_var\_usage}(\overline{s})$}
    \AxiomC{$\Gamma' \tack \overline{s} \tackl \Gamma''$}
    \RightLabel{[T-Constructor]}
    \TrinaryInfC{$\Gamma \tack \mathtt{new}\ k(\overline{x : \phi_p})\ \{\ \overline{s}\ \} \tackl \Gamma$}
\end{prooftree}

\vspace{1em}

\begin{prooftree}
    \AxiomC{$\Gamma' = \Gamma[V \mapsto \{\overline{x : \phi_p}\}]$}
    \AxiomC{$\Gamma \tack \text{sendable}(\overline{\phi_p})$}
    \AxiomC{$\Gamma' \tack \text{valid\_var\_usage}(\overline{s})$}
    \AxiomC{$\Gamma' \tack \overline{s} \tackl \Gamma''$}
    \RightLabel{[T-Be]}
    \QuaternaryInfC{$\Gamma \tack \mathtt{be}\ b(\overline{x : \phi_p})\ \{\ \overline{s}\ \} \tackl \Gamma$}
\end{prooftree}

\vspace{1em}

\begin{prooftree}
    \AxiomC{$\Gamma \tack \overline{T} : \text{valid}$}
    \AxiomC{$\Gamma_{\text{base}} = \Gamma[this \mapsto \text{Actor}(A)]$}
    \AxiomC{$\Gamma_{\text{base}} \tack \overline{M} \tackl \Gamma_{\text{base}}$}
    \RightLabel{[T-Actor]}
    \TrinaryInfC{$\Gamma \tack \mathtt{actor}\ A\ \{\ \overline{f : T};\ \overline{M}\ \} \tackl \Gamma$}
\end{prooftree}
\end{center}

\chapter{\textit{UserMutex implementation}} \label{appendix:user_mutex_implementation}
\section{Locking}
\begin{algorithmic}[1]
\Procedure{Lock}{$locker\_id$}
    \State \textbf{acquire} $coord\_lock$
    
    % Path 1: Unowned
    \If{$holding\_instance = \text{null}$}
        \State $holding\_instance \gets locker\_id$
        \State $mutex.num\_lock\_called \gets 1$
        \State \textbf{release} $mutex.coord\_lock$
        \State \Return \textbf{true}
    \EndIf

    % Path 2: Reentrant
    \If{$holding\_instance = locker\_id$}
        \State $num\_lock\_called \gets num\_lock\_called + 1$
        \State \textbf{release} $mutex.coord\_lock$
        \State \Return \textbf{true}
    \EndIf

    % Path 3: Contention
    \State Set $locker\_id$'s state to \textbf{WAITING} and add it to mutex's $wait\_queue$
    \State \textbf{release} $mutex.coord\_lock$
    \State \Return \textbf{false}
\EndProcedure
\end{algorithmic}
\section{Unlocking}
\begin{algorithmic}[1]
\Procedure{Unlock}{$runtime$}
    \State \textbf{acquire} $mutex.coord\_lock$
    \State $mutex.num\_lock\_called \gets mutex.num\_lock\_called - 1$
    
    \If{$mutex.num\_lock\_called > 0$}
        \State \textbf{release} $mutex.coord\_lock$
        \State \Return
    \EndIf

    \If{$mutex.wait\_queue$ is empty}
        \State $mutex.holding\_instance \gets \text{null}$
    \Else
        \State $target \gets mutex.wait\_queue.pop\_front()$
        \State $mutex.holding\_instance \gets target$
        \State $mutex.num\_lock\_called \gets 1$
        
        \State Set $target$'s state to \textbf{RUNNABLE} and add it to the \textit{schedule\_queue}
    \EndIf
    \State \textbf{release} $mutex.coord\_lock$
\EndProcedure
\end{algorithmic}

\chapter{Model checker} \label{appendix:prolog}
\begin{lstlisting}[language = prolog]
% The maximum occurences of reference capabilities in the 
% predicates below is 3, and hence only 3 different locks are 
% needed.
lock(a).
lock(b).
lock(c).

% All valid capabilities
capability(iso).
capability(iso_cap).
capability(val).
capability(ref).
capability(tag).
capability(locked(X)) :- lock(X).

% view(K1, K2, K3) succeeds if K1 > K2 = K3, 
% where > is the viewpoint 
adaptation operator
view(iso, ref, iso).
view(iso, val, val).
view(iso, tag, tag).
view(iso, iso, iso).
view(iso_cap, ref, iso_cap).    
view(iso_cap, val, val).    
view(iso_cap, tag, tag).
view(iso_cap, iso, iso_cap).
view(val, ref, val).
view(val, val, val).
view(val, tag, tag).
view(val, iso, val).
view(ref, ref, ref).
view(ref, val, val).
view(ref, tag, tag).
view(ref, iso, iso).
view(tag, ref, tag).
view(tag, val, tag).
view(tag, iso, tag).
view(tag, tag, tag).
view(locked(X), tag, tag) :- lock(X).
view(locked(X), ref, locked(X)) :- lock(X).
view(locked(X), val, val) :- lock(X).
view(locked(X), iso, iso) :- lock(X).
view(locked(X), locked(Y), locked(Y)) :- lock(X), lock(Y).

% check_nonassociativity(K1, K2, K3): succeeds if (K1 > K2) > K3 is
% not equal to (K1 > (K2 > K3)).
check_nonassociativity(K1, K2, K3) :- 
    view(K1, K2, K1_2), view(K1_2, K3, K1_2_3), 
    view(K2, K3, K2_3), \+ view(K1, K2_3, K1_2_3).

% covariant(K1, K2): succeeds if K1 is covariant to K2
covariant(tag, X) :- capability(X).
covariant(X, X) :- capability(X), \+ (X = iso_cap).
covariant(X, iso_cap) :- capability(X), \+ (X = iso_cap).

% invariant(K1, K2): succeeds if K1 is invariant to K2
invariant(X, X) :- capability(X), \+ (X = iso_cap).
invariant(X, iso_cap) :- capability(X), \+ (X = iso_cap).

% assignable(K1, K2): succeeds if K2 is assignable to K1
assignable(tag, X) :- capability(X).
assignable(X, X) :- capability(X), \+ (member(X, [iso, iso_cap])).
assignable(X, iso_cap) :- capability(X), \+ (X = iso_cap).

% mutable(K): succeeds if K is a mutable capability
mutable(K) :- 
    capability(K), member(K, [ref, iso, iso_cap, locked(L)]).

% readable(K): succeeds if K can be read but not written
readable(val).

% opaque(K): succeeds if K can neither be written to nor read 
% from
opaque(tag).

% ptr_assignable(K, K1, K2): succeeds if the type (T K) K2 is 
% assignable to (T K) K1
ptr_assignable(K, K1, K2) :- assignable(K1, K2),
    view(K1, K, ViewK1), view(K2, K, ViewK2),
    mutable(K1), invariant(ViewK1, ViewK2).
ptr_assignable(K, K1, K2) :- 
    assignable(K1, K2), view(K1, K, ViewK1), view(K2, K, ViewK2), 
    readable(K1), covariant(ViewK1, ViewK2).
ptr_assignable(K, K1, K2) :- 
    assignable(K1, K2), capability(K), capability(K2), opaque(K1).

% check_nested_assignment_inconsistencies(K, K1, K2): succeeds 
% if K2 is assignable to K1, but the pointer type (T K) K2 is 
% not assignable to (T K) K1.
check_nested_assignment_inconsistencies(K, K1, K2) :- 
    assignable(K1, K2), \+ ptr_assignable(K, K1, K2).

% ptr_covariant(K, K1, K2): succeeds if the type (T K) K1 is 
% covariant to (T K) K2
ptr_covariant(K, K1, K2) :- 
    covariant(K1, K2), view(K1, K, ViewK1), view(K2, K, ViewK2), 
    mutable(K1), invariant(ViewK1, ViewK2).
ptr_covariant(K, K1, K2) :- 
    covariant(K1, K2), view(K1, K, ViewK1), view(K2, K, ViewK2), 
    readable(K1), covariant(ViewK1, ViewK2).
ptr_covariant(K, K1, K2) :- 
    covariant(K1, K2), capability(K), capability(K2), opaque(K1).

% check_nested_covariant_inconsistencies(K, K1, K2): succeeds if 
% K1 is covariant to K2, but the pointer type (T K) K1 is not 
% covariant to (T K) K2.
check_nested_covariant_inconsistencies(K, K1, K2) :- 
    covariant(K1, K2), \+ ptr_covariant(K, K1, K2).

% ptr_invariant(K, K1, K2): succeeds if the type (T K) K1 is 
% invariant to (T K) K2
ptr_invariant(K, K1, K2) :- 
    invariant(K1, K2), capability(K), capability(K2), opaque(K1).
ptr_invariant(K, K1, K2) :- 
    invariant(K1, K2), view(K1, K, ViewK1), 
    view(K2, K, ViewK2), invariant(ViewK1, ViewK2).

% check_nested_invariant_inconsistencies(K, K1, K2): succeeds if 
% K1 is invariant to K2, but the pointer type (T K) K1 is not 
% invariant to (T K) K2.
check_nested_invariant_inconsistencies(K, K1, K2) :- 
    invariant(K1, K2), \+ ptr_invariant(K, K1, K2).

% check_iso_unalias_inconsistencies(K): succeeds if the type 
% iso > K is not compatible with iso_cap > K.
check_iso_unalias_inconsistencies(K) :- 
    view(iso, K, iso), view(iso_cap, K, iso_cap), 
    \+ (invariant(IsoK, IsocapK)).
check_iso_unalias_inconsistencies(K) :-
    view(iso, K, IsoK), view(iso_cap, K, iso_cap), 
    \+ (covariant(IsocapK, IsoK)).


% sendable(K): succeeds if K is sendable
sendable(K) :- 
    capability(K), member(K, [iso, val, locked(L), tag]).
% iso_holdable(K): succeeds if having an iso reference as a 
% member of a pointer type protected by K does not break the
% alias invariant of iso. The idea is that only synchronized 
% types, or types guaranteeing the single alias invariant are 
% allowed to hold and send an iso.
iso_holdable(K) :- capability(K), member(K, [iso, locked(L)]).

% check_sendable_inconsistencies(K1, K2) succeeds if the rule 
% for K1 > K2 is leads to unsoundness when a pointer of type (T 
% K2) K1 is present as a parameter of a behaviour. 
check_sendable_inconsistencies(K1, K2) :- 
    sendable(K1), view(K1, K2, ref).
check_sendable_inconsistencies(K1, K2) :- 
    sendable(K1), \+ (iso_holdable(K1)), view(K1, K2, iso).

% Checks all the inconsistencies

check_inconsistencies :- check_nonassociativity(K1, K2, K3).
check_inconsistencies :- 
    check_nested_assignment_inconsistencies(K, K1, K2).
check_inconsistencies :- 
    check_nested_covariant_inconsistencies(K, K1, K2).
check_inconsistencies :- 
    check_nested_invariant_inconsistencies(K, K1, K2).
check_inconsistencies :- check_iso_unalias_inconsistencies(K).
check_inconsistencies :- check_sendable_inconsistencies(K1, K2).
\end{lstlisting}

\chapter{Hot-potato pattern in Pony} \label{appendix:hot-potato}
The following code implements a scheduler that manages access to a shared integer. Actors request access by calling \textsf{take} with a callback. When the resource is available, the callback is invoked with exclusive ownership of the resource. After completing their operation, actors return the resource via \textsf{give}.
\begin{lstlisting}
class IntContainer
  var value: U64 = 0
  new create(v: U64) => value = v

actor IntScheduler
  var _data: (IntContainer iso | None) = None
  var _wait_queue: Array[{(IntContainer iso): None} val]
   = Array[{(IntContainer iso): None} val]

  be take(cb: {(IntContainer iso): None} val) =>
    match _data = None
    | let p: IntContainer iso => cb(consume p)
    | None => _wait_queue.push(cb)
    end

  be give(p: IntContainer iso) =>
    if _wait_queue.size() > 0 then
      try
        let cb = _wait_queue.shift()?
        cb(consume p)
      end
    else
      _data = consume p
    end
\end{lstlisting}

The following code demonstrates how the hot-potato pattern can be used to atomically swap the values of two shared integers. Note that the operation must be split across three behaviours (\textsf{swap}, \textsf{receive1}, \textsf{receive2}) because Pony behaviours cannot block. Intermediate state is stored in the actor field \textsf{\_p1}.
\begin{lstlisting}
actor Swapper
  var _s1: IntScheduler
  var _s2: IntScheduler
  var _p1: (IntContainer iso | None) = None

  new create(s1: IntScheduler, s2: IntScheduler) =>
    _s1 = s1
    _s2 = s2

  be receive2(p2: IntContainer iso) =>
    // Step 3: Perform the swap
    match _p1 = None
    | let p1: IntContainer iso =>
      let temp = p1.value
      p1.value = p2.value
      p2.value = temp
      // Step 4: Return both
      _s1.give(consume p1)
      _s2.give(consume p2)
    end

  be receive1(p1: IntContainer iso) =>
    _p1 = consume p1
    let self: Swapper tag = this
    _s2.take({(p2: IntContainer iso): None =>
      self.receive2(consume p2)
      None} val)

  be swap() =>
    let self: Swapper tag = this
    _s1.take({(p1: IntContainer iso): None =>
      self.receive1(consume p1)
      None} val)
\end{lstlisting}
